\chapter{因子消融与边界扰动}
\label{ch:ablate}

主对照把食品名、厚度、初温、含水、非均质和策略一次全部打开。要判断哪一个因子真正移动停机偏差，必须先固定其余因子，再单独打开一条。厨房里对应的问题更具体：换一块土豆，停机符号会不会翻转？从冷藏改到室温，会不会少走一段焓升？表面结了一层霜，跟踪器会不会把开局读数当成已经热了？下面先读 S1 对偶子集上的边际，再读 S2 在 30 mm 肉糜上的边界扫描；两条都出自同一数值核。含水三档与质量通道开、关给出零效应，脂肪夹层和开门给出可测的负向移动，遮挡时长把表面读数保持在初值若干秒之后，中位从 \(-0.112\) 抬到 \(-0.067\)。

\section{食品属性的边际贡献}

对偶策略在 S1 中有 4860 条。按食品名拆开，肉糜、土豆、面团的中位 \(e_{\alpha}\) 分别是 \(-0.338\)、\(-0.325\)、\(-0.313\)，十分位都在 \(-0.767\)。三种名称几乎重合，因为停止判据读取的是冷点进度，活度只调制界面通量的幅度；GAB 清洁子集里肉类条数为 0，核里的活度是含水的单调饱和函数。面板上的“鱼、蔬菜、馒头”若只改名称、不改厚度和初温，交叉表给出的中位几乎不动。名称容易进入菜单文案，真正移动停机偏差的是厚度和脂肪夹层。

按初始含水拆开，\(m_0\in\{0.64,0.72,0.80\}\) 三条分位表完全相同：中位 \(-0.313\)，十分位 \(-0.767\)，九十分位 \(0.000\)。这是一次真零效应。界面潜热\eqref{eq:qlat}在活度被裁到上沿之后对含水不再敏感，结壳又要求表面含水低于 0.45，而 S1 的初始含水达不到这条门槛。含水作为因子进入格子，格子却把它读成了常数。

按初温拆开，\(5^\circ\mathrm{C}\) 的中位是 \(-0.418\)，\(20^\circ\mathrm{C}\) 是 \(-0.183\)。从冷藏到室温少走大约 \(15\,\mathrm{K}\) 的内部升温，停机偏差上移 \(0.23\)。内部焓决定还要走完的温升：探针窗在两种初温上都同样执行，差别落在初值上。

按非均质拆开，均匀场中位 \(-0.246\)，一层夹心 \(-0.158\)，两层 \(-0.439\)。两层把导热在指定控制体上乘 0.45，冷点更难起来；一层夹心的中位反而高于均匀场，因为夹心位置不必落在最低温度格子上，部分种子的冷点被推到更薄的一侧。图\ref{fig:s1het}给出这三档的分位。

\begin{figure}[htbp]
\centering
\includegraphics[width=0.88\textwidth]{X_s1_hetero.png}
\caption{S1 对偶策略：非均质档对停机偏差的分位。两层夹心把中位压到 \(-0.439\)。}
\label{fig:s1het}
\end{figure}

按厚度拆开已经在第\ref{ch:result}章给出：15 mm 中位 \(-0.417\)，30 mm \(-0.080\)，50 mm \(-0.598\)。厚度的极差大于食品名、含水和初温。图\ref{fig:s1fac}把四个因子的边际画在同一张图上，厚度轴的跨度最大，含水轴是一条水平线。

\begin{figure}[htbp]
\centering
\includegraphics[width=0.88\textwidth]{X_s1_factors.png}
\caption{S1 对偶策略的因子边际。厚度主导；初始含水三条重合。}
\label{fig:s1fac}
\end{figure}

\section{边界与结构扰动的影响}

S2 固定在 30 mm、肉糜、初温 \(5^\circ\mathrm{C}\)、含水 0.72，合计 480 行，用来看边界扰动。表\ref{tab:s2}给出各组中位。

\begin{table}[htbp]
\centering
\caption{S2 边界扫描的停机偏差中位。各组 \(n=40\)。}
\label{tab:s2}
\begin{tabular}{llrr}
\toprule
因子 & 水平 & \(e_{50}\) & 十分位--九十分位 \\
\midrule
探针 & 对偶 & \(-0.112\) & \([-0.155,-0.111]\) \\
探针 & 无探针 & \(-0.579\) & \([-0.579,-0.578]\) \\
读数保持 & \(0\,\mathrm{s}\) & \(-0.112\) & \([-0.155,-0.111]\) \\
读数保持 & \(15\,\mathrm{s}\) & \(-0.111\) & 四十条重合 \\
读数保持 & \(30\,\mathrm{s}\) & \(-0.110\) & 四十条重合 \\
读数保持 & \(60\,\mathrm{s}\) & \(-0.067\) & 四十条重合 \\
质量通道 & 开 & \(-0.112\) & 与关重合 \\
质量通道 & 关 & \(-0.112\) & 与开重合 \\
脂肪夹层 & 关 & \(-0.112\) & \([-0.155,-0.111]\) \\
脂肪夹层 & 开 & \(-0.353\) & \([-0.354,-0.313]\) \\
门 & 关闭 & \(-0.112\) & \([-0.155,-0.111]\) \\
门 & 300 s 开门 & \(-0.151\) & \([-0.151,-0.107]\) \\
\bottomrule
\end{tabular}
\end{table}

对偶相对无探针把中位从 \(-0.579\) 抬到 \(-0.112\)。无探针四十条几乎没有散布，跟踪器以同一滞后追上同一表面；对偶在窗后接受热阻，停止时刻随种子有几十秒的移动，十分位落到 \(-0.155\)。图\ref{fig:s2probe}给出两组的分位。

\begin{figure}[htbp]
\centering
\includegraphics[width=0.88\textwidth]{X_s2_probe.png}
\caption{S2：对偶与无探针的停机偏差。无探针四十条几乎重合在 \(-0.579\)。}
\label{fig:s2probe}
\end{figure}

开局遮挡把表面观测保持在初值，时长取 \(0/15/30/60\,\mathrm{s}\)。核在 \(t<\texttt{occlude\_s}\) 时用开局读数替换当前表面温度，到达该时长之后才放开跟踪。从 \(0\,\mathrm{s}\) 到 \(30\,\mathrm{s}\)，中位几乎不动；\(60\,\mathrm{s}\) 上移到 \(-0.067\)。更长的保持让表面读数偏低，跟踪器更晚宣布到达，真值进度多走了一段：观测偏差改写了停止时刻，导热本身没有变。核没有热像通道，这四档代替的是开局若干秒内表面读数被固定在初值。图\ref{fig:s2occ}给出四档，图\ref{fig:d15occ}把同一扫描画成中位随保持秒数的折线：60 s 仍在零下。

\begin{figure}[htbp]
\centering
\includegraphics[width=0.88\textwidth]{X_s2_occlude.png}
\caption{S2：开局保持表面读数的时长。\(60\,\mathrm{s}\) 让停止推迟，真值进度上移。}
\label{fig:s2occ}
\end{figure}

\begin{figure}[htbp]
\centering
\includegraphics[width=0.88\textwidth]{D15_occlusion.png}
\caption{S2：表面读数保持初值的秒数对照中位停机偏差。60 s 仍在零下。核没有热像通道。}
\label{fig:d15occ}
\end{figure}

质量通道打开与关闭给出完全相同的分位。跟踪器并不读取质量，S2 把这一开关放进格子，核把它读成了空操作。厨房里称重托盘能看见累计蒸发，看不见冷点位置；这条零效应与表\ref{tab:obs}里质量通道的物理内容一致。脂肪夹层打开后中位从 \(-0.112\) 降到 \(-0.353\)：指定控制体导热乘 0.45，冷点被拖住。一块带肥膘的肉糜，导热在夹层处下降，菜单时间仍按“肉类、中等”给出，停机偏差被夹层单独压下去。300 s 开门一次，中位从 \(-0.112\) 降到 \(-0.151\)，腔温回落把表面通量打断几十秒。厨房里对应的是中途开门查看、或者往里续一盘：门微动开关切断加热和蒸汽，冷点在这几十秒里几乎不动，表面却先凉一截。图\ref{fig:s2three}把质量、脂肪、开门三组画在一起：质量两条重合，脂肪是本组里移动最大的边界。

\begin{figure}[htbp]
\centering
\includegraphics[width=0.88\textwidth]{X_s2_three.png}
\caption{S2：质量通道、脂肪夹层与 300 s 开门。质量开与关重合；脂肪夹层把中位压到 \(-0.353\)。}
\label{fig:s2three}
\end{figure}

图\ref{fig:s2radar}把各因子相对对偶基线的中位移画成横条。无探针把中位压下约 0.47，脂肪夹层压下约 0.24，300 s 开门压下约 0.04，开局保持 60 s 把中位抬高约 0.045，质量通道关几乎贴着零。横条只用于同图比较，不定义新的综合分数。图\ref{fig:s2fac}给出 S2 全因子的箱线。固定条件是 30 mm、肉糜、初温 \(5^\circ\mathrm{C}\)、含水 0.72。在已经看见热阻的 30 mm 上，脂肪和开门把中位拉低，遮挡把中位抬高，质量通道保持为零。厚度这一主导因素仍由第\ref{ch:result}章的交叉表给出。

\begin{figure}[htbp]
\centering
\includegraphics[width=0.86\textwidth]{S2_factors.png}
\caption{S2 全因子箱线。脂肪夹层与无探针离开基线最远；质量开与关重合。}
\label{fig:s2fac}
\end{figure}

\begin{figure}[htbp]
\centering
\includegraphics[width=0.86\textwidth]{X_s2_radar.png}
\caption{S2 各因子相对对偶基线的中位移。无探针与脂肪夹层离开基线最远。}
\label{fig:s2radar}
\end{figure}

\section{含水变化的轨迹复核}

S1 的含水零效应在单条轨迹上同样成立。种子 7、30 mm、组合档、对偶，\(m_0=0.64\) 时停机偏差 \(+0.143\)、停止 1856 s；\(m_0=0.72\) 与 \(0.80\) 都是 \(+0.145\)、1866 s。三条结束方差都是 0.224。含水改变界面通量的幅度，不足以在这一条轨迹上改写停止时刻的符号。

固定 30 mm、肉糜之后，对偶相对无探针把中位抬了约 \(0.47\)，脂肪夹层把同一位置压下约 \(0.24\)，开门一次压下约 \(0.04\)；开局保持表面读数 60 s 让真值进度多走。质量通道与初始含水按零效应报告，食品名称的边际小于厚度和非均质。这些句子与 S1 全表的混合中位分开读：先按厚度拆开，再按脂肪和探针拆开。

面板上的食品名最容易被当成决定因素。鱼、蔬菜、馒头、肉类四个字对应四条菜单时间，说明书也把它们印成四列，每一列给出一个分钟数。S1 对偶子集里肉糜、土豆、面团的中位分别是 \(-0.338\)、\(-0.325\)、\(-0.313\)，三种原型只差 0.025，十分位都在 \(-0.767\)。名称改的是活度尺度，停机读取的是冷点进度，活度只调制界面通量的幅度，因此名称几乎不动停机偏差。含水三档 \(0.64/0.72/0.80\) 给出同一条分位，因为界面潜热在活度被裁到上沿之后对含水不再敏感，结壳又要求表面含水低于 0.45，而 S1 的初始含水达不到这条门槛。厨房里称重托盘能看见累计蒸发，看不见冷点在哪一层；S2 把质量通道开与关读成重合的两条分位。含水和质量通道都按零效应报告。

真正移动中位的是脂肪夹层、中途开门和有没有那 40 s 探针窗。带肥膘的肉糜把指定控制体导热乘 0.45，冷点被拖住，S2 中位从 \(-0.112\) 降到 \(-0.353\)，移动幅度大于开门、小于无探针相对对偶的约 0.47。中途开门在厨房里几乎不可避免：看一眼上色、续一盘、擦一下门玻璃，门微动落下，加热和蒸汽被置为零，腔温回落把表面通量打断几十秒；S2 把这一次记作 300 s 开门，中位从 \(-0.112\) 降到 \(-0.151\)。演示页上的开门按钮只改铰链开度，不改已经写出的温度数组，为的是让零件可见，而不是在浏览器里重算开门扰动；真值扰动记录在 S2 的 480 行里。菜单仍按“肉类、中等”给出分钟数时，夹层、开门和探针窗这三件事都不在文案里。固定 30 mm 之后再读 S2，才能把折进菜单表的那三件事重新拆开。

开局结霜是说明书几乎不写、厨房里却天天遇到的一件事。从冷冻室取出的肉糜表面结了一层薄霜，红外或门玻璃内侧的表面通道会把读数保持在冰点附近几十秒；S2 把这一段记作 60 s 霜层，中位从 \(-0.112\) 抬到 \(-0.067\)：跟踪器更晚宣布到达，真值进度反而多走了一段。观测偏差改写了停止时刻，导热本身没有变。核把 \(-7.92^\circ\mathrm{C}\) 取为先验热阻 5 的滞后上限，演示页 HUD 标成开局余量。固定 30 mm 之后，能改停机符号的边界是探针、脂肪和开门，不是食品名和含水。食品名留在菜单表上，积分器读取的是厚度、夹层和那 40 s 窗口。
